\documentclass[12pt,a4paper]{report}
\usepackage[utf8]{inputenc}
\usepackage[T1]{fontenc}
\usepackage[french]{babel}
\usepackage{geometry,setspace,graphicx,float,xcolor,booktabs,longtable,array,hyperref,enumitem,listings,tikz}
\usetikzlibrary{positioning,arrows.meta,shapes.geometric}
\geometry{margin=2.5cm}
\onehalfspacing
\graphicspath{{figures/}}
\setlength{\tabcolsep}{4pt}
\setlength{\LTleft}{0pt}
\setlength{\LTright}{0pt}
\setlength{\emergencystretch}{3em}
\hypersetup{colorlinks=true,linkcolor=black,urlcolor=blue,citecolor=black}
\definecolor{brandolive}{RGB}{90,90,64}
\definecolor{codebg}{RGB}{248,248,248}
\definecolor{codeframe}{RGB}{210,210,210}
\lstdefinelanguage{JavaScript}{keywords={async,await,const,let,var,for,of,if,else,return,new,true,false,null,undefined,try,catch,import,from,export,function},sensitive=true,comment=[l]{//},morecomment=[s]{/*}{*/},morestring=[b]',morestring=[b]",morestring=[b]`}
\lstdefinestyle{codestyle}{backgroundcolor=\color{codebg},frame=single,rulecolor=\color{codeframe},basicstyle=\ttfamily\small,breaklines=true,columns=fullflexible,showstringspaces=false,keywordstyle=\color{blue},commentstyle=\color{gray},stringstyle=\color{teal}}
\lstset{style=codestyle}
\newcommand{\placeholder}[1]{\textbf{[À compléter : #1]}}
\newcommand{\screenshotplaceholder}[3]{\begin{figure}[H]\centering\fbox{\begin{minipage}[c][6.2cm][c]{0.88\textwidth}\centering\textbf{Capture à insérer}\par\vspace{0.3cm}#1\par\vspace{0.3cm}Fichier recommandé : \texttt{docs/rapport/figures/#2}\end{minipage}}\caption{#3}\label{fig:#2}\end{figure}}
\newcommand{\reportfigure}[3]{\begin{figure}[p]\centering\includegraphics[width=\textwidth,height=0.72\textheight,keepaspectratio]{#1}\caption{#2}\label{fig:#3}\end{figure}}
\begin{document}
\begin{titlepage}
\centering
\vspace*{1cm}
{\Large \placeholder{nom de l'institution}\par}
{\large \placeholder{département ou faculté}\par}
\vspace{1.8cm}
{\Huge\bfseries Saveur : application de recettes de cuisine\par}
\vspace{0.4cm}
{\Large Rapport académique de projet d'examen\par}
\vspace{1.3cm}
\begin{tabular}{rl}
\textbf{Projet :} & Projet 5 -- Application de Recettes de Cuisine \\
\textbf{Cours :} & Développement JavaScript \\
\textbf{Cours :} & Bases de données NoSQL \\
\textbf{Étudiants :} & \placeholder{noms et prénoms des étudiants du groupe} \\
\textbf{Groupe :} & \placeholder{numéro ou nom du groupe} \\
\textbf{Enseignant :} & \placeholder{nom de l'enseignant} \\
\textbf{Date de remise :} & \placeholder{date de remise officielle} \\
\textbf{Année académique :} & 2025--2026 \\
\end{tabular}
\vfill
{\large Rapport rédigé à partir de l'audit du dépôt Saveur\par}
\end{titlepage}

\chapter*{Résumé}
\addcontentsline{toc}{chapter}{Résumé}
Saveur est une application web full-stack de recettes de cuisine réalisée dans le cadre des cours de Développement JavaScript et de Bases de données NoSQL. Nous avons conçu une plateforme permettant de consulter un catalogue public de recettes, de filtrer les résultats par catégorie, temps et difficulté, de créer des recettes personnelles, de gérer des favoris, de planifier les repas de la semaine et de générer une liste de courses intelligente. Le frontend utilise React, TypeScript, Vite, Tailwind CSS et des composants dynamiques. Le backend repose sur Node.js, Express, MongoDB et Mongoose. L'authentification est assurée par un couple email/mot de passe, un hash bcrypt et un token JWT envoyé aux routes protégées. La base contient sept collections MongoDB liées par références et sous-documents. Les vérifications disponibles confirment que le typage TypeScript, les tests automatisés de la liste de courses et le build de production passent localement. Certaines validations restent manuelles, notamment les captures finales et les tests navigateur complets.

\textbf{Mots-clés :} JavaScript, React, Express, MongoDB, Mongoose, JWT, recettes, liste de courses.
\tableofcontents
\listoffigures
\listoftables

\chapter{Introduction}
Le présent projet s'inscrit dans le cadre des cours de Développement JavaScript et de Bases de données NoSQL. Ces deux enseignements visent à consolider les compétences nécessaires à la création d'une application web complète : interface dynamique côté client, serveur applicatif Node.js, API REST, persistance NoSQL et mécanismes de sécurité. Le document de projet exige explicitement une application frontend et backend, l'utilisation de JavaScript ES6+, Node.js/Express, MongoDB avec Mongoose, une authentification complète et une organisation claire selon les responsabilités du modèle, de la vue et du contrôleur.

Le sujet attribué à notre groupe est le Projet 5 : application de recettes de cuisine. Ce domaine implique une modélisation plus riche qu'un simple catalogue : une recette contient des informations descriptives, des ingrédients avec quantités et unités, des étapes de préparation, une catégorie, une cuisine, une image et un auteur. L'application doit aussi permettre à l'utilisateur de planifier ses repas et de transformer ce planning en liste de courses. Cette transformation est un élément central du sujet, car elle demande de parcourir plusieurs documents, de regrouper des ingrédients compatibles et de conserver l'état des articles déjà cochés.

Nous avons développé Saveur comme application web utilisable localement. L'utilisateur visite une page d'accueil orientée découverte, parcourt le catalogue, consulte les détails d'une recette, se connecte, crée ses propres recettes, ajoute des favoris, planifie ses repas et génère ses courses. Le tableau de bord personnel est placé dans la page de profil, où l'utilisateur retrouve ses indicateurs, son carnet de recettes, ses favoris, ses préférences de langue et les liens vers le planning.

Ce rapport suit la structure officielle demandée par le modèle de rapport. Il présente l'analyse des besoins, la conception technique et NoSQL, l'implémentation, la sécurité, les fonctionnalités, les tests, les difficultés rencontrées, les améliorations possibles et les annexes nécessaires pour la remise académique.

\chapter{Analyse des besoins}
\section{Présentation du problème}
La gestion de recettes devient rapidement complexe lorsque l'on ajoute les contraintes de la vie quotidienne. Un utilisateur peut avoir plusieurs recettes, différentes unités de mesure, des ingrédients répétés, des favoris, des repas à organiser dans la semaine et une liste de courses à préparer. Sans outil centralisé, les informations restent dispersées : les recettes sont consultées à un endroit, le planning est écrit ailleurs, et la liste de courses doit être reconstruite manuellement.

Saveur répond à ce problème en regroupant ces opérations dans une même application. Les recettes servent de source de données principale. Le planning hebdomadaire réutilise ces recettes. La liste de courses est générée à partir des ingrédients des repas planifiés. Les favoris et le carnet personnel permettent à l'utilisateur de retrouver rapidement les contenus importants. Le catalogue public facilite la découverte de recettes existantes et répond à l'exigence académique de filtrage.

\section{Objectifs fonctionnels}
Le tableau \ref{tab:functional} présente les besoins fonctionnels, leur preuve dans le dépôt et leur état actuel. Les statuts distinguent les éléments implémentés, partiels, manquants ou non vérifiés.

\begin{longtable}{p{2.5cm}p{3.7cm}p{2.3cm}p{3.9cm}p{2cm}}
\caption{Tableau des exigences fonctionnelles}\label{tab:functional}\\
\toprule
\textbf{Exigence} & \textbf{Description} & \textbf{Utilisateur} & \textbf{Preuve dans le dépôt} & \textbf{Statut} \\
\midrule
\endfirsthead
\toprule
\textbf{Exigence} & \textbf{Description} & \textbf{Utilisateur} & \textbf{Preuve dans le dépôt} & \textbf{Statut} \\
\midrule
\endhead
Authentification & Inscription, connexion, session utilisateur & Visiteur, utilisateur & \texttt{server/controllers/authController.js}, \texttt{server/middleware/auth.js} & Implémenté \\
Mot de passe sécurisé & Hashage avant stockage & Utilisateur & \texttt{bcrypt.hash(password, 12)} & Implémenté \\
Recettes de cuisine & Création et lecture de recettes avec description, temps, catégorie et image & Utilisateur, visiteur & \texttt{server/models/Recipe.js}, \texttt{src/pages/CreateRecipe.tsx} & Implémenté \\
Ingrédients avec unités & Chaque ingrédient de recette possède quantité et unité & Utilisateur & \texttt{Recipe.ingredients[]} & Implémenté \\
Étapes de préparation & Étapes ordonnées avec titre et description & Utilisateur, visiteur & \texttt{Recipe.steps[]} & Implémenté \\
Catalogue public & Liste des recettes publiques & Visiteur, utilisateur & \texttt{GET /api/recipes}, \texttt{Catalog.tsx} & Implémenté \\
Filtres catalogue & Catégorie, temps maximal, difficulté et recherche & Visiteur, utilisateur & \texttt{listRecipes}, \texttt{getFilteredRecipes} & Implémenté \\
Favoris & Ajouter ou retirer une recette favorite & Utilisateur & \texttt{favoriteController.js} & Implémenté \\
Recettes personnelles & Lister les recettes créées par l'utilisateur & Utilisateur & \texttt{GET /api/recipes/mine} & Implémenté \\
Tableau de bord & Statistiques personnelles dans le profil & Utilisateur & \texttt{dashboardController.js}, \texttt{Profile.tsx} & Implémenté \\
Planning hebdomadaire & Planifier repas par date et type & Utilisateur & \texttt{MealPlan.js}, \texttt{Plan.tsx} & Implémenté \\
Liste de courses intelligente & Génération depuis le planning, agrégation et catégories de magasin & Utilisateur & \texttt{groceryListService.js}, \texttt{shoppingCategories.js} & Implémenté \\
Téléversement d'image & Upload local, capture caméra, URL externe & Utilisateur & \texttt{uploadRoutes.js}, \texttt{CreateRecipe.tsx} & Implémenté localement \\
Internationalisation & Interface en anglais, français, espagnol, italien & Tous & \texttt{src/services/i18n.tsx} & Partiel à vérifier \\
Mode PWA & Manifest, service worker, page hors ligne & Tous & \texttt{public/manifest.webmanifest}, \texttt{public/sw.js} & Base implémentée \\
Édition de recette & Modifier une recette après création & Utilisateur & Aucun endpoint \texttt{PUT/PATCH /api/recipes/:id} trouvé & Manquant \\
Administration visible & Interface administrateur & Admin & Rôle \texttt{admin} dans \texttt{User.js}, pas d'écran admin vérifié & Partiel \\
\bottomrule
\end{longtable}

\section{Objectifs non fonctionnels}
Sur le plan de la sécurité, Saveur utilise bcrypt pour éviter le stockage des mots de passe en clair et JWT pour protéger les routes sensibles. Le middleware \texttt{requireAuth} vérifie le jeton et charge l'utilisateur avant d'autoriser l'accès. Certaines protections d'appartenance existent : la suppression d'une recette vérifie que l'utilisateur est propriétaire ou administrateur, et les ressources de planning, favoris, dashboard et courses sont filtrées par \texttt{req.user.\_id}.

La maintenabilité est favorisée par une organisation claire du code. Le backend sépare les modèles Mongoose, routes Express, contrôleurs, middleware, services et utilitaires. Le frontend sépare les pages React, services API, types TypeScript, composants partagés et utilitaires. Cette structure facilite l'audit et répond à la consigne de séparation des responsabilités.

La performance est traitée de manière simple mais suffisante pour un MVP académique. Les listes de recettes sont limitées à 100 éléments côté backend. Les filtres sont appliqués en base pour la catégorie, la difficulté, la recherche textuelle et le temps maximal. Le catalogue applique ensuite une pagination côté interface. Un index texte existe sur le titre et la description des recettes.

L'ergonomie a fait l'objet de corrections successives : navigation desktop sans débordement, page d'accueil séparée du dashboard, profil comme espace personnel, catalogue paginé, liste de courses regroupée par sections de magasin et formulaire de recette avec options d'image. L'accessibilité n'a pas été auditée par un outil spécialisé ; nous la classons donc comme partiellement vérifiée. La fiabilité repose sur la gestion des erreurs Express et sur des messages JSON lorsque les requêtes échouent. Le service worker ajoute une base PWA, mais l'édition hors ligne et la synchronisation différée ne sont pas implémentées.

\section{Utilisateurs et rôles}
Le visiteur non connecté peut consulter les pages publiques, notamment l'accueil, le catalogue et les détails des recettes publiques. Il ne peut pas créer de recette, ajouter de favori, planifier de repas ou générer une liste de courses. Les actions protégées redirigent vers la page de profil/connexion ou sont bloquées par l'API.

L'utilisateur authentifié peut créer des recettes, consulter son carnet, gérer ses favoris, remplir le planning, générer la liste de courses, cocher les articles et accéder au dashboard personnel. Les données personnelles sont liées à son identifiant MongoDB.

Le rôle administrateur existe dans le modèle \texttt{User} et dans le token JWT. Le backend autorise déjà la suppression d'une recette par un administrateur. Toutefois, aucune interface d'administration complète n'a été vérifiée dans le dépôt. Nous considérons donc l'administration comme une possibilité technique partielle, non comme une fonctionnalité complète.

\section{Cas d'utilisation principaux}
Les principaux parcours sont l'inscription et la connexion, la consultation du catalogue, le filtrage des recettes, la création d'une recette personnelle, l'ajout aux favoris, la planification d'un repas, la génération de la liste de courses et la gestion du profil. La création de recette enregistre les ingrédients avec quantité et unité. La génération des courses lit les recettes planifiées sur la semaine, agrège les ingrédients compatibles, conserve les articles personnalisés et présente les résultats par catégories de magasin.

\chapter{Conception technique}
\section{Architecture logicielle}
L'application suit une architecture trois tiers. La couche cliente est une application React/Vite. Elle consomme une API REST exposée par Express. La couche de données est MongoDB, manipulée via Mongoose. En production locale, Express peut aussi servir le dossier \texttt{dist} généré par Vite.

\begin{figure}[H]
\centering
\begin{tikzpicture}[node distance=2cm, box/.style={rectangle, rounded corners, draw=brandolive, thick, align=center, minimum width=4cm, minimum height=1.1cm}, arrow/.style={-{Latex}, thick}]
\node[box] (client) {Client React/Vite\\Pages, services, i18n};
\node[box, right=of client] (api) {API Express\\Routes, contrôleurs, middleware};
\node[box, right=of api] (db) {MongoDB\\Collections Mongoose};
\node[box, below=of api] (uploads) {Stockage local\\uploads/recipes};
\draw[arrow] (client) -- node[above]{HTTP JSON / FormData} (api);
\draw[arrow] (api) -- node[above]{Mongoose} (db);
\draw[arrow] (api) -- (uploads);
\draw[arrow] (api) to[bend left=20] node[below]{JWT Bearer} (client);
\end{tikzpicture}
\caption{Architecture logicielle vérifiée de Saveur}
\label{fig:architecture}
\end{figure}

\section{Organisation générale du système}
Le frontend porte la responsabilité de l'affichage, de la navigation, de la gestion locale de la session, de la traduction des libellés et des interactions utilisateur. Les appels API sont regroupés dans \texttt{src/services}. Les pages principales sont \texttt{Home.tsx}, \texttt{Catalog.tsx}, \texttt{RecipeDetail.tsx}, \texttt{CreateRecipe.tsx}, \texttt{Plan.tsx}, \texttt{Groceries.tsx} et \texttt{Profile.tsx}.

Le backend centralise les règles métier. Les routes sont montées dans \texttt{server/index.js} sous les préfixes \texttt{/api/auth}, \texttt{/api/categories}, \texttt{/api/ingredients}, \texttt{/api/recipes}, \texttt{/api/favorites}, \texttt{/api/meal-plans}, \texttt{/api/groceries}, \texttt{/api/dashboard} et \texttt{/api/uploads}. Les contrôleurs manipulent les modèles Mongoose et renvoient des réponses JSON. La base conserve les comptes, catégories, ingrédients canoniques, recettes, favoris, entrées de planning et listes de courses. Les images téléversées sont stockées localement dans \texttt{uploads/recipes}.

\section{Modélisation NoSQL}
Le modèle NoSQL combine références et sous-documents. Les références permettent de relier les utilisateurs, catégories, recettes, favoris et plans. Les sous-documents conviennent aux ingrédients d'une recette, aux étapes et aux articles de courses, car ces données sont généralement chargées avec leur document parent.

\begin{longtable}{p{2.2cm}p{3.2cm}p{4.2cm}p{4.8cm}}
\caption{Détail des collections MongoDB}\label{tab:collections}\\
\toprule
\textbf{Collection} & \textbf{Rôle} & \textbf{Champs principaux} & \textbf{Contraintes et relations} \\
\midrule
\endfirsthead
\toprule
\textbf{Collection} & \textbf{Rôle} & \textbf{Champs principaux} & \textbf{Contraintes et relations} \\
\midrule
\endhead
\texttt{users} & Comptes utilisateur & \texttt{name}, \texttt{email}, \texttt{passwordHash}, \texttt{role}, \texttt{photoURL}, \texttt{preferredLanguage} & Email unique, rôle \texttt{user/admin}, langue \texttt{en/fr/es/it}. Référencé par recettes, favoris, plans et listes. \\
\texttt{categories} & Collections de recettes & \texttt{name}, \texttt{slug}, \texttt{image} & Nom unique, slug obligatoire. Référencé par \texttt{recipes}. \\
\texttt{ingredients} & Catalogue canonique d'ingrédients & \texttt{name}, \texttt{defaultUnit}, \texttt{category}, \texttt{caloriesPerUnit} & Nom unique, catégorie de courses stable, unité par défaut. \\
\texttt{recipes} & Recettes et contenu culinaire & \texttt{title}, \texttt{description}, temps, difficulté, portions, cuisine, image, visibilité, \texttt{ingredients[]}, \texttt{steps[]} & Référence \texttt{user}, référence \texttt{category}, ingrédients et étapes embarqués, index texte titre/description. \\
\texttt{favorites} & Favoris utilisateur & \texttt{user}, \texttt{recipe} & Index unique \texttt{user + recipe}. Relation de jonction. \\
\texttt{mealplans} & Planning hebdomadaire & \texttt{user}, \texttt{date}, \texttt{mealType}, \texttt{recipe} & Date au format ISO, type de repas enum, index unique \texttt{user + date + mealType}. \\
\texttt{shoppinglists} & Courses par semaine & \texttt{user}, \texttt{weekStart}, \texttt{items[]} & Index unique \texttt{user + weekStart}. Articles embarqués : nom, quantité, unité, catégorie, source, checked. \\
\bottomrule
\end{longtable}

\section{Relations entre les collections}
Un utilisateur peut créer plusieurs recettes. Une catégorie peut classer plusieurs recettes. Les favoris constituent une relation plusieurs-à-plusieurs entre utilisateurs et recettes. La collection \texttt{favorites} joue le rôle de document de jonction, avec un index unique pour éviter qu'un utilisateur ajoute deux fois la même recette. Le planning relie un utilisateur à une recette pour une date et un type de repas. La liste de courses est associée à l'utilisateur et à un début de semaine. Ses articles sont embarqués, car ils représentent le résultat calculé et modifiable d'une génération.

\begin{figure}[H]
\centering
\begin{tikzpicture}[node distance=1.35cm and 1.45cm, col/.style={rectangle, rounded corners, draw=brandolive, fill=brandolive!8, thick, align=center, minimum width=2.75cm, minimum height=0.9cm}, rel/.style={-{Latex}, thick, draw=brandolive}]
\node[col] (users) {users};
\node[col, right=of users] (recipes) {recipes};
\node[col, above=of recipes] (categories) {categories};
\node[col, below=of recipes] (ingredients) {ingredients};
\node[col, right=of recipes] (favorites) {favorites};
\node[col, below=of favorites] (mealplans) {mealplans};
\node[col, below=of mealplans] (shopping) {shoppinglists};
\draw[rel] (users) -- node[above]{1 crée N} (recipes);
\draw[rel] (categories) -- node[right]{1 classe N} (recipes);
\draw[rel] (ingredients) -- node[below]{1 apparaît dans N} (recipes);
\draw[rel] (users) -- node[above,sloped]{1 possède N} (favorites);
\draw[rel] (recipes) -- node[above,sloped]{1 reçoit N} (favorites);
\draw[rel] (users) -- node[below,sloped]{1 possède N} (mealplans);
\draw[rel] (recipes) -- node[below,sloped]{1 planifiée N} (mealplans);
\draw[rel] (users) to[bend right=18] node[left]{1 possède N} (shopping);
\draw[rel,dashed] (mealplans) -- node[right]{alimente} (shopping);
\end{tikzpicture}
\caption{Diagramme des collections MongoDB et de leurs relations principales}
\label{fig:relations}
\end{figure}

\section{Justification des choix NoSQL}
MongoDB convient à ce projet parce qu'une recette est naturellement un document riche. Elle contient des métadonnées, une catégorie, une visibilité, des ingrédients et des étapes. Les ingrédients et les étapes n'ont pas besoin d'être interrogés indépendamment dans la majorité des cas ; leur embarquement réduit donc le nombre de requêtes lors de l'affichage d'une fiche recette. Les références restent nécessaires pour les entités réutilisables. Les utilisateurs, catégories, recettes planifiées et favoris sont reliés par identifiants. La collection \texttt{ingredients} apporte une couche de normalisation utile pour la liste de courses, car chaque ingrédient canonique possède une unité par défaut et une catégorie de magasin stable.

\section{Authentification et autorisation}
Lors de l'inscription, le client envoie nom, email et mot de passe à \texttt{POST /api/auth/register}. Le contrôleur vérifie la présence des champs et impose un mot de passe d'au moins six caractères. Il recherche ensuite un éventuel compte existant par email. Si l'email est libre, il appelle \texttt{bcrypt.hash(password, 12)} et crée l'utilisateur avec le hash. Le serveur retourne un JWT et un objet utilisateur sérialisé.

Lors de la connexion, \texttt{POST /api/auth/login} vérifie l'email et compare le mot de passe saisi avec \texttt{bcrypt.compare}. Si la comparaison réussit, \texttt{createToken} crée un JWT contenant \texttt{sub} et \texttt{role}, avec une expiration de sept jours. Côté frontend, \texttt{AuthContext.tsx} stocke la session dans \texttt{localStorage} sous la clé \texttt{saveur\_session}. Le service \texttt{apiFetch} lit ensuite cette session et ajoute l'en-tête \texttt{Authorization: Bearer <token>}.

Le middleware \texttt{requireAuth} extrait le token, vérifie sa signature avec \texttt{JWT\_SECRET}, charge l'utilisateur sans \texttt{passwordHash} et place l'objet dans \texttt{req.user}. Les routes protégées peuvent alors utiliser \texttt{req.user.\_id} pour limiter les données à l'utilisateur connecté.

\begin{longtable}{p{1.4cm}p{3.5cm}p{4.4cm}p{2cm}p{3.1cm}}
\caption{Inventaire des routes API importantes}\label{tab:routes}\\
\toprule
\textbf{Méthode} & \textbf{Route} & \textbf{Rôle} & \textbf{Auth} & \textbf{Fichier} \\
\midrule
\endfirsthead
\toprule
\textbf{Méthode} & \textbf{Route} & \textbf{Rôle} & \textbf{Auth} & \textbf{Fichier} \\
\midrule
\endhead
POST & \texttt{/api/auth/register} & Créer un compte & Non & \texttt{authRoutes.js} \\
POST & \texttt{/api/auth/login} & Connecter un utilisateur & Non & \texttt{authRoutes.js} \\
GET & \texttt{/api/auth/me} & Récupérer la session & Oui & \texttt{authRoutes.js} \\
GET & \texttt{/api/categories} & Lister les catégories & Non & \texttt{categoryRoutes.js} \\
GET & \texttt{/api/ingredients} & Chercher les ingrédients & Non & \texttt{ingredientRoutes.js} \\
GET & \texttt{/api/recipes} & Catalogue public filtré & Optionnelle & \texttt{recipeRoutes.js} \\
GET & \texttt{/api/recipes/:id} & Détail recette & Optionnelle & \texttt{recipeRoutes.js} \\
POST & \texttt{/api/recipes} & Créer une recette & Oui & \texttt{recipeRoutes.js} \\
GET & \texttt{/api/recipes/mine} & Recettes personnelles & Oui & \texttt{recipeRoutes.js} \\
DELETE & \texttt{/api/recipes/:id} & Supprimer une recette & Oui & \texttt{recipeRoutes.js} \\
GET & \texttt{/api/favorites} & Lister les favoris & Oui & \texttt{favoriteRoutes.js} \\
POST & \texttt{/api/favorites/:recipeId/toggle} & Ajouter/retirer favori & Oui & \texttt{favoriteRoutes.js} \\
GET & \texttt{/api/meal-plans} & Lire le planning & Oui & \texttt{mealPlanRoutes.js} \\
POST & \texttt{/api/meal-plans} & Ajouter/remplacer repas & Oui & \texttt{mealPlanRoutes.js} \\
GET & \texttt{/api/groceries} & Générer/lire les courses & Oui & \texttt{groceryRoutes.js} \\
POST & \texttt{/api/groceries/custom} & Ajouter article personnalisé & Oui & \texttt{groceryRoutes.js} \\
PATCH & \texttt{/api/groceries/:itemId} & Cocher/décocher article & Oui & \texttt{groceryRoutes.js} \\
GET & \texttt{/api/dashboard} & Statistiques profil & Oui & \texttt{dashboardRoutes.js} \\
POST & \texttt{/api/uploads/recipe-image} & Téléverser image recette & Oui & \texttt{uploadRoutes.js} \\
\bottomrule
\end{longtable}

\chapter{Implémentation}
\section{Stack technologique}
\begin{longtable}{p{2.8cm}p{2.2cm}p{4.4cm}p{5cm}}
\caption{Technologies utilisées}\label{tab:stack}\\
\toprule
\textbf{Technologie} & \textbf{Version} & \textbf{Rôle} & \textbf{Justification} \\
\midrule
\endfirsthead
\toprule
\textbf{Technologie} & \textbf{Version} & \textbf{Rôle} & \textbf{Justification} \\
\midrule
\endhead
React & \texttt{\^{}19.0.1} & Interface utilisateur & Composants dynamiques et SPA. \\
TypeScript & \texttt{\~{}5.8.2} & Typage frontend & Limiter les erreurs de structure de données. \\
Vite & \texttt{\^{}6.2.3} & Build et serveur dev & Démarrage rapide et build optimisé. \\
Tailwind CSS & \texttt{\^{}4.1.14} & Styles & Interface responsive et cohérente. \\
Node.js & 24.14.0 observé & Runtime backend & Exécution JavaScript serveur. \\
Express & \texttt{\^{}4.21.2} & API REST & Routage HTTP et middleware. \\
MongoDB & Local & Base NoSQL & Documents flexibles et relations par références. \\
Mongoose & \texttt{\^{}8.9.5} & ODM & Schémas, validations et requêtes MongoDB. \\
bcryptjs & \texttt{\^{}2.4.3} & Hashage & Sécurisation des mots de passe. \\
jsonwebtoken & \texttt{\^{}9.0.2} & JWT & Sessions stateless côté API. \\
multer & \texttt{\^{}2.2.0} & Upload & Réception des images de recette. \\
Vitest & \texttt{\^{}4.1.5} & Tests & Tests unitaires du service de courses. \\
\bottomrule
\end{longtable}

\section{Organisation du code}
\begin{lstlisting}[language=bash,caption={Arborescence commentee du projet}]
server/
  config/db.js                   # connexion Mongoose
  controllers/                   # logique des endpoints REST
  data/seedData.js               # categories, ingredients, recettes, comptes demo
  middleware/auth.js             # JWT, optionalAuth, requireRole
  models/                        # schemas User, Recipe, Category...
  routes/                        # definition des routes Express
  services/groceryListService.js # generation et aggregation des courses
  utils/shoppingCategories.js    # categories et alias d'ingredients
src/
  components/                    # layout, selecteur de langue
  pages/                         # Home, Catalog, Plan, Groceries, Profile...
  services/                      # appels API, auth, i18n
  utils/                         # icones de categories, categories de courses
  types.ts                       # types partages frontend
public/
  manifest.webmanifest           # PWA manifest
  sw.js                          # service worker
  offline.html                   # fallback hors ligne
\end{lstlisting}

\section{Fonctionnalités principales implémentées}
\subsection{Authentification}
L'authentification est implémentée côté backend avec \texttt{authController.js}. L'inscription crée un utilisateur après vérification de l'email et hashage du mot de passe. La connexion compare le mot de passe avec le hash existant. Le frontend expose le formulaire de connexion et d'inscription dans \texttt{Profile.tsx}. Les champs de démonstration sont préremplis avec le compte \texttt{chef@saveur.local}, ce qui facilite les tests locaux. Cette donnée est explicitement présente dans le seed et le README ; elle doit être changée ou retirée dans un environnement de production.

\subsection{Recettes, ingrédients et étapes}
La création de recette est disponible pour les utilisateurs connectés. Le formulaire contient les informations générales, la catégorie, la cuisine, l'image, les temps, le niveau de difficulté, les portions, les ingrédients et les étapes. Côté serveur, \texttt{createRecipe} vérifie que la catégorie existe, filtre les étapes vides et normalise les ingrédients. Une recette privée ne peut être consultée que par son propriétaire ou par un administrateur. La suppression exige que l'utilisateur soit propriétaire de la recette ou administrateur. En revanche, l'édition complète d'une recette existante n'a pas été vérifiée dans le dépôt.

\subsection{Catalogue, favoris, profil et planning}
Le catalogue utilise \texttt{GET /api/recipes}. Les filtres vérifiés sont la catégorie, la difficulté, le temps maximal et la recherche texte. Les favoris sont stockés dans une collection dédiée, ce qui évite d'embarquer une longue liste dans le document utilisateur. Le profil charge les recettes personnelles, les favoris et les statistiques du dashboard. Le planning est stocké dans \texttt{mealplans}; chaque document contient l'utilisateur, la date, le type de repas et la recette. Un index unique empêche deux recettes d'occuper le même créneau pour le même utilisateur.

\subsection{Liste de courses intelligente}
Le contrôleur \texttt{generateGroceryList} récupère les repas planifiés entre \texttt{weekStart} et \texttt{weekStart + 6 jours}. Il collecte les identifiants des ingrédients canoniques, charge leurs catégories, récupère l'ancienne liste pour conserver les éléments cochés et appelle \texttt{buildGeneratedGroceryItems}. Les articles personnalisés existants sont ensuite réintégrés.

Les catégories de magasin sont stables : \texttt{DAIRY}, \texttt{PRODUCE}, \texttt{PASTA\_GRAINS}, \texttt{PANTRY}, \texttt{MEAT\_PROTEIN}, \texttt{BAKERY}, \texttt{FROZEN}, \texttt{BEVERAGES} et \texttt{OTHER}. Les libellés visibles sont traduits côté frontend. Des alias permettent de reconnaître des variantes comme \texttt{Greek yogurt}, \texttt{yaourt grec}, \texttt{tomatoes}, \texttt{tagliatelle} ou \texttt{walnuts}. Lorsqu'aucune correspondance fiable n'est trouvée, l'article tombe dans \texttt{OTHER}.

\subsection{Images, langues et PWA}
Le formulaire conserve trois options d'image : URL externe, upload local et capture caméra. L'endpoint \texttt{POST /api/uploads/recipe-image} accepte un champ multipart \texttt{photo}, limite la taille à 5 Mo et autorise JPG, PNG et WebP. L'interface contient des dictionnaires pour l'anglais, le français, l'espagnol et l'italien. Le projet contient aussi un manifest, une icône, un service worker et une page \texttt{offline.html}. Nous considérons donc la PWA comme une base installable, non comme une application offline-first complète.

\section{Extraits de code commentés}
\begin{lstlisting}[language=JavaScript,caption={Middleware d'authentification -- server/middleware/auth.js}]
const header = req.headers.authorization || '';
const token = header.startsWith('Bearer ') ? header.slice(7) : null;
if (!token) {
  return res.status(401).json({ message: 'Authentication required.' });
}
const payload = jwt.verify(token, process.env.JWT_SECRET);
const user = await User.findById(payload.sub).select('-passwordHash');
\end{lstlisting}
Cet extrait montre comment les routes privées refusent les requêtes non authentifiées, vérifient le token et chargent l'utilisateur sans exposer le hash du mot de passe.

\begin{lstlisting}[language=JavaScript,caption={Normalisation d'ingrédients -- server/controllers/recipeController.js}]
const ingredient = await Ingredient.findOneAndUpdate(
  { name: new RegExp(`^${name.replace(/[.*+?^${}()|[\]\\]/g, '\\$&')}$`, 'i') },
  {
    $setOnInsert: { name, defaultUnit: unit },
    $set: { category: suggestShoppingCategory(name, item.category) },
  },
  { new: true, upsert: true }
);
\end{lstlisting}
Cet extrait prouve que les ingrédients ne sont pas seulement manipulés dans l'interface. Le backend crée ou met à jour un ingrédient canonique, en tenant compte d'une suggestion de catégorie de courses.

\begin{lstlisting}[language=JavaScript,caption={Agrégation des courses -- server/services/groceryListService.js}]
const key = groceryItemKey({ name: recipeIngredient.name, unit });
const previous = previousByKey.get(key);
const current = itemMap.get(key) || {
  name: recipeIngredient.name,
  quantity: 0,
  unit,
  category: suggestShoppingCategory(recipeIngredient.name, canonicalIngredient?.category),
  checked: Boolean(previous?.checked),
  source: 'generated',
};
current.quantity += Number(recipeIngredient.quantity) || 0;
itemMap.set(key, current);
\end{lstlisting}
Cet extrait montre la logique essentielle : normaliser la clé d'article, conserver l'état coché précédent, résoudre la catégorie et additionner les quantités compatibles.

\section{Interface utilisateur et expérience utilisateur}
La direction visuelle utilise des tons olive et neutres, des titres serif, des cartes arrondies et des ombres discrètes. La navigation principale contient Accueil, Catalogue, Planning et Courses. Le profil est accessible par l'avatar, car il contient le dashboard personnel. La page d'accueil sert à la découverte : recette du jour, catégories et recommandations. Le catalogue propose les collections, les filtres et la pagination. Le profil regroupe l'identité utilisateur, les statistiques, le carnet, les favoris et les préférences. Le planning se concentre sur la semaine et les types de repas. La liste de courses suit une organisation pratique par section de magasin. Le détail recette sépare les ingrédients et les étapes, avec un mode cuisine focalisé.

\chapter{Tests et validation}
\section{Stratégie et environnement de test}
La stratégie de test combine contrôles automatisés et validation manuelle. Les contrôles automatisés présents couvrent le service de liste de courses : reconnaissance d'alias, agrégation de quantités, conservation de l'état coché et ajout personnalisé. Le typage TypeScript sert de contrôle statique via \texttt{npm run lint}. Le build de production vérifie que l'application peut être compilée avec Vite. Aucun test end-to-end navigateur complet n'a été trouvé dans le dépôt.

Les commandes ont été exécutées dans le dossier \texttt{C:\\Users\\migue\\JS\_NoSQL\\Saveur---Recipe-Meal-Planner}. L'environnement observé utilise Windows, Node.js 24.14.0 dans le terminal utilisateur, MongoDB local sur \texttt{127.0.0.1:27017}, Vite pour le frontend et Express pour le backend.

\section{Plan de tests}
\begin{longtable}{p{0.9cm}p{2.1cm}p{2.5cm}p{3.9cm}p{2.8cm}p{1.5cm}}
\caption{Plan de tests}\label{tab:tests}\\
\toprule
\textbf{ID} & \textbf{Fonction} & \textbf{Préconditions} & \textbf{Étapes} & \textbf{Résultat attendu} & \textbf{Statut} \\
\midrule
\endfirsthead
\toprule
\textbf{ID} & \textbf{Fonction} & \textbf{Préconditions} & \textbf{Étapes} & \textbf{Résultat attendu} & \textbf{Statut} \\
\midrule
\endhead
T01 & Inscription & API active & Créer un compte valide & Token et utilisateur retournés & À valider \\
T02 & Connexion demo & Seed actif & Utiliser \texttt{chef@saveur.local} & Session créée & À valider \\
T03 & Mauvais mot de passe & Compte existant & Mot de passe incorrect & Erreur 401 & À valider \\
T04 & Route protégée & Pas de token & Appeler \texttt{/api/auth/me} & Erreur 401 & À valider \\
T05 & Catalogue & Seed actif & Ouvrir \texttt{/catalog} & Recettes affichées & À valider \\
T06 & Filtre catégorie & Catalogue chargé & Choisir une collection & Résultats filtrés & À valider \\
T07 & Filtre temps & Catalogue chargé & Choisir un temps maximal & Recettes sous limite & À valider \\
T08 & Filtre difficulté & Catalogue chargé & Choisir une difficulté & Recettes filtrées & À valider \\
T09 & Création recette & Connecté & Remplir formulaire & Recette créée & À valider \\
T10 & Favori & Connecté & Cliquer sur coeur & Favori modifié & À valider \\
T11 & Planning & Connecté & Choisir repas et recette & Entrée créée & À valider \\
T12 & Courses agrégation & Tests unitaires & Exécuter \texttt{npm run test} & Quantités agrégées & OK \\
T13 & Catégorie courses & Tests unitaires & Exécuter \texttt{npm run test} & Alias reconnus & OK \\
T14 & État coché & Tests unitaires & Régénérer une liste simulée & \texttt{checked} conservé & OK \\
T15 & Build production & Dépendances installées & Exécuter \texttt{npm run build} & \texttt{dist} généré & OK \\
T16 & Responsive & Navigateur & Tester mobile/tablette/desktop & Pas de débordement & À valider \\
\bottomrule
\end{longtable}

\section{Résultats des tests}
Les commandes suivantes ont été exécutées avec succès le 12 juillet 2026 : \texttt{npm.cmd run lint} a terminé \texttt{tsc --noEmit} sans erreur ; \texttt{npm.cmd run test} a exécuté 1 fichier de test avec 4 tests réussis ; \texttt{npm.cmd run build} a transformé 2102 modules et généré le dossier \texttt{dist}. Ces résultats confirment que le typage, le service de courses et la compilation de production sont opérationnels.

\chapter{Difficultés rencontrées et solutions}
\section{Adaptation du projet au cahier des charges web}
Le projet avait initialement une orientation mobile. Le sujet de l'enseignant impose une application web avec frontend, backend Node.js/Express et MongoDB. Nous avons donc recentré l'application sur React/Vite côté client, Express côté serveur et MongoDB/Mongoose côté données. Le résultat est une structure full-stack cohérente avec les exigences, mais les captures finales doivent encore remplacer les anciennes captures annotées.

\section{Modélisation des ingrédients}
La génération d'une liste de courses exige de connaître les quantités et les unités. Une simple liste de noms d'ingrédients aurait été insuffisante. La solution retenue consiste à embarquer dans chaque recette les ingrédients avec \texttt{name}, \texttt{quantity}, \texttt{unit} et \texttt{optional}, tout en gardant une collection \texttt{ingredients} réutilisable. La limite restante est l'absence de conversion entre unités différentes.

\section{Catégorisation intelligente de la liste de courses}
Une liste de courses utile doit être organisée selon les sections du magasin. La catégorie ne devait pas être décidée uniquement dans le composant React. Nous avons donc ajouté des catégories stables côté domaine, une table d'alias et une fonction \texttt{suggestShoppingCategory}. Les ingrédients et articles de courses conservent une catégorie interne. Les inconnus restent classés dans \texttt{OTHER} jusqu'à correction manuelle ou enrichissement du catalogue.

\section{Protection des ressources personnelles}
Les recettes personnelles, favoris, plans et listes de courses ne doivent pas être mélangés entre utilisateurs. Les routes protégées utilisent \texttt{requireAuth}. Les requêtes MongoDB filtrent généralement par \texttt{req.user.\_id}. La suppression de recette vérifie la propriété ou le rôle administrateur. Un audit de sécurité complet reste nécessaire avant production ; l'utilisation de \texttt{localStorage} pour le JWT est acceptable pour le projet, mais des cookies HTTP-only seraient préférables en production.

\section{Cohérence de l'interface et navigation responsive}
Les corrections récentes ont porté sur le navbar desktop, la page d'accueil, le profil et la liste de courses. L'application avait évolué rapidement depuis un prototype ; certaines responsabilités d'écran étaient mélangées. Le dashboard a été déplacé dans le profil, l'accueil a été orienté découverte, la navigation a été simplifiée et la liste de courses a été restructurée en sections. Les captures finales doivent encore confirmer visuellement chaque écran dans l'état de soumission.

\chapter{Améliorations possibles}
\begin{longtable}{p{2.2cm}p{3.8cm}p{4.6cm}p{3.8cm}}
\caption{Améliorations futures}\label{tab:improvements}\\
\toprule
\textbf{Priorité} & \textbf{Limitation actuelle} & \textbf{Solution proposée} & \textbf{Bénéfice attendu} \\
\midrule
\endfirsthead
\toprule
\textbf{Priorité} & \textbf{Limitation actuelle} & \textbf{Solution proposée} & \textbf{Bénéfice attendu} \\
\midrule
\endhead
Court terme & Peu de tests automatisés & Ajouter tests API et composants React & Réduire les régressions \\
Court terme & Édition de recette non vérifiée & Ajouter endpoint \texttt{PATCH /api/recipes/:id} et écran d'édition & Gestion complète des recettes \\
Moyen terme & Stockage image local & Migrer vers Cloudinary, S3 ou équivalent & Persistance après déploiement \\
Moyen terme & JWT stocké dans \texttt{localStorage} & Utiliser cookies HTTP-only et durcissement CSRF & Sécurité renforcée \\
Moyen terme & Internationalisation à valider & Relire tous les libellés et caractères & Interface plus professionnelle \\
Long terme & Pas de conversion d'unités & Ajouter normalisation g/kg/ml/l & Courses plus exactes \\
Long terme & Pas d'avis utilisateurs & Ajouter notation, commentaires et modération & Découverte enrichie \\
Long terme & PWA limitée & Ajouter cache de données et synchronisation différée & Hors ligne plus complet \\
\bottomrule
\end{longtable}

\chapter{Conclusion}
Nous avons atteint l'objectif principal : produire une application web full-stack de recettes conforme à l'esprit du sujet. Saveur contient un frontend dynamique, un backend Express, une base MongoDB avec sept collections, une authentification par JWT et bcrypt, des recettes avec ingrédients et étapes, un catalogue filtrable, des favoris, un profil/dashboard, un planning hebdomadaire et une liste de courses intelligente.

Le projet nous a permis de renforcer nos compétences en JavaScript moderne, en architecture REST, en modélisation NoSQL et en séparation des responsabilités. La mise en place de Mongoose nous a aidés à structurer les documents tout en profitant de la flexibilité MongoDB. L'authentification nous a permis de manipuler les notions de hashage, token, middleware et routes protégées. Côté frontend, React et TypeScript ont facilité l'organisation en pages, services et composants.

L'application n'est pas encore un produit final complet. Les tests navigateur doivent être documentés, les captures finales doivent être produites, l'édition des recettes reste à ajouter, et les améliorations de sécurité seront nécessaires avant une mise en production réelle. Malgré ces limites, Saveur constitue une base fonctionnelle et défendable pour la remise académique et la soutenance.

\begin{thebibliography}{9}
\addcontentsline{toc}{chapter}{Références}
\bibitem{brief} Document d'examen : \textit{PROJETS D'EXAMEN DE DEVELOPPEMENT JAVASCRIPT \& BASE DE DONNEES NOSQL}, fourni par l'enseignant.
\bibitem{model} Modèle de rapport : \textit{MODEL DU RAPPORT PROJETS D'EXAMEN DE DEVELOPPEMENT JAVASCRIPT \& BASE DE DONNEES NOSQL}, fourni par l'enseignant.
\bibitem{node} Node.js Documentation, \url{https://nodejs.org/docs/latest/api/}.
\bibitem{express} Express Documentation, \url{https://expressjs.com/}.
\bibitem{mongodb} MongoDB Manual, \url{https://www.mongodb.com/docs/}.
\bibitem{mongoose} Mongoose Documentation, \url{https://mongoosejs.com/docs/}.
\bibitem{react} React Documentation, \url{https://react.dev/}.
\bibitem{vite} Vite Documentation, \url{https://vite.dev/}.
\bibitem{jwt} JSON Web Tokens, \url{https://jwt.io/introduction}.
\end{thebibliography}

\appendix
\chapter{Captures d'écran}
Le modèle de rapport exige au moins six captures. Les figures A.1 à A.8 présentent les principaux écrans vérifiés de l'application dans leur état final. Elles montrent respectivement l'accueil, l'authentification, le catalogue, le profil, le planning, la liste de courses, le formulaire de création et la fiche détaillée d'une recette.
\reportfigure{home.png}{Page d'accueil de Saveur avec recette mise en avant et catégories}{annex-home}
\reportfigure{auth.png}{Interface d'authentification utilisateur}{annex-auth}
\reportfigure{catalog.png}{Catalogue de recettes avec collections et filtres}{annex-catalog}
\reportfigure{profile-dashboard.png}{Profil et tableau de bord de l'utilisateur authentifié}{annex-profile}
\reportfigure{planner.png}{Planning hebdomadaire des repas}{annex-planner}
\reportfigure{groceries.png}{Liste de courses intelligente regroupée par rayons}{annex-groceries}
\reportfigure{create-recipe.png}{Formulaire de création d'une recette avec image, ingrédients et étapes}{annex-create-recipe}
\reportfigure{recipe-detail.png}{Détail d'une recette avec ingrédients et accès au mode cuisine}{annex-recipe-detail}

\chapter{Guide d'installation et d'exécution}
\section*{Prérequis}
\begin{itemize}[leftmargin=*]
\item Node.js installé.
\item MongoDB Server installé et démarré localement.
\item Git pour récupérer le dépôt.
\item Un terminal PowerShell, Git Bash ou équivalent.
\end{itemize}
\section*{Étapes}
\begin{enumerate}[leftmargin=*]
\item Installer les dépendances : \texttt{npm install}.
\item Copier \texttt{.env.example} vers \texttt{.env}.
\item Configurer \texttt{PORT}, \texttt{MONGO\_URI}, \texttt{JWT\_SECRET}, \texttt{CLIENT\_URL}, \texttt{AUTO\_SEED} et \texttt{VITE\_API\_URL}.
\item Vérifier que MongoDB fonctionne sur \texttt{127.0.0.1:27017}.
\item Lancer l'API avec \texttt{npm run api}.
\item Dans un autre terminal, lancer le frontend avec \texttt{npm run dev}.
\item Ouvrir \texttt{http://localhost:3000}. Si ce port est occupé, Vite peut proposer un autre port.
\item Vérifier le build avec \texttt{npm run build} et les tests avec \texttt{npm run test}.
\end{enumerate}
\section*{Compte de démonstration}
Le seed local crée un compte de démonstration destiné aux tests : \texttt{chef@saveur.local} avec le mot de passe \texttt{saveur123}. Ce compte ne doit pas être utilisé comme identifiant réel en production.

\chapter{Principales routes API}
Le tableau \ref{tab:routes} dans le chapitre de conception contient l'inventaire principal des routes. Pour la soutenance, nous recommandons de démontrer au minimum : connexion, liste des recettes, création de recette, ajout au planning, génération de courses et récupération du dashboard.

\chapter{Répartition des tâches}
Les informations de groupe ne sont pas disponibles dans le dépôt. Le tableau suivant doit être complété avant la remise.
\begin{table}[H]
\centering
\caption{Répartition des tâches du groupe}
\begin{tabular}{p{3.2cm}p{4.5cm}p{3.5cm}p{3.2cm}}
\toprule
\textbf{Étudiant} & \textbf{Contributions principales} & \textbf{Contributions secondaires} & \textbf{Tests/documentation} \\
\midrule
\placeholder{étudiant 1} & Backend, modèles, routes & Aide frontend & Tests API \\
\placeholder{étudiant 2} & Interface React, UX & Traductions & Captures \\
\placeholder{étudiant 3} & MongoDB, seed, liste de courses & Planning & Rapport \\
\placeholder{étudiant 4 si applicable} & Intégration & Corrections & Relecture \\
\bottomrule
\end{tabular}
\end{table}

\chapter{Table de traçabilité interne}
\begin{longtable}{p{4.5cm}p{6.3cm}p{3.6cm}}
\caption{Correspondance entre affirmations du rapport et preuves du dépôt}\label{tab:evidence}\\
\toprule
\textbf{Affirmation} & \textbf{Preuve dépôt} & \textbf{Statut} \\
\midrule
\endfirsthead
\toprule
\textbf{Affirmation} & \textbf{Preuve dépôt} & \textbf{Statut} \\
\midrule
\endhead
Authentification JWT & \texttt{server/utils/tokens.js}, \texttt{server/middleware/auth.js} & Vérifié \\
Hashage bcrypt & \texttt{server/controllers/authController.js} & Vérifié \\
Sept collections MongoDB & \texttt{server/models/*.js} & Vérifié \\
Recettes avec ingrédients et étapes & \texttt{server/models/Recipe.js} & Vérifié \\
Catalogue filtrable & \texttt{recipeController.js}, \texttt{Catalog.tsx} & Vérifié \\
Planning hebdomadaire & \texttt{MealPlan.js}, \texttt{Plan.tsx} & Vérifié \\
Liste de courses intelligente & \texttt{groceryListService.js}, \texttt{Groceries.tsx} & Vérifié \\
PWA de base & \texttt{public/manifest.webmanifest}, \texttt{public/sw.js} & Vérifié \\
Tests automatisés & \texttt{server/services/groceryListService.test.js} & Vérifié \\
Édition complète des recettes & Aucun endpoint d'édition trouvé & Manquant \\
\bottomrule
\end{longtable}
\end{document}
